\documentclass[12pt,a4paper]{article}

% ============================================================
% PACKAGES
% ============================================================
\usepackage[utf8]{inputenc}
\usepackage[T1]{fontenc}
\usepackage[vietnamese]{babel}
\usepackage{lmodern}
\usepackage[top=2.5cm, bottom=2.5cm, left=2.5cm, right=2.5cm]{geometry}
\usepackage{amsmath}
\usepackage{graphicx}
\usepackage{tikz}
\usetikzlibrary{shapes.geometric, arrows.meta, positioning, fit, calc, decorations.pathreplacing}
\usepackage{booktabs}
\usepackage{hyperref}
\usepackage{enumitem}
\usepackage{xcolor}
\usepackage{setspace}
\usepackage{fancyhdr}
\usepackage{mdframed}

\onehalfspacing

% Mau sac
\definecolor{core}{HTML}{2E86AB}
\definecolor{ext1}{HTML}{A23B72}
\definecolor{ext2}{HTML}{F18F01}
\definecolor{opt}{HTML}{6B8E23}
\definecolor{lightgray}{HTML}{F5F5F5}
\definecolor{accepted}{HTML}{27AE60}
\definecolor{rejected}{HTML}{E74C3C}

% Header/Footer
\pagestyle{fancy}
\fancyhf{}
\fancyhead[L]{\small\textit{Giải thích Pipeline -- Hệ thống nhận dạng từ khóa thoại}}
\fancyhead[R]{\small\thepage}
\renewcommand{\headrulewidth}{0.4pt}

\hypersetup{colorlinks=true, linkcolor=core, urlcolor=core}

% Box cho vi du
\newmdenv[
  backgroundcolor=lightgray,
  linecolor=core,
  linewidth=1pt,
  roundcorner=5pt,
  innertopmargin=8pt,
  innerbottommargin=8pt,
  innerleftmargin=10pt,
  innerrightmargin=10pt,
  skipabove=10pt,
  skipbelow=10pt
]{examplebox}

% ============================================================
\title{
    \vspace{-1cm}
    {\Large\textbf{Hệ Thống Nhận Dạng Từ Khóa Thoại}} \\[0.3cm]
    {\large Giải thích toàn bộ quy trình} \\[0.2cm]
    {\normalsize\textit{Dành cho người không có nền tảng kỹ thuật}}
}
\author{}
\date{}

\begin{document}
\maketitle
\thispagestyle{fancy}

\tableofcontents
\newpage


% ============================================================
\section{Dự án này làm gì?}
% ============================================================

Bạn đã bao giờ nói ``Hey Siri'' với iPhone, hoặc ``OK Google'' với điện thoại Android
chưa? Khi bạn nói câu đó, điện thoại \textbf{nhận ra từ khóa} trong giọng nói của bạn
và ``thức dậy'' để lắng nghe tiếp. Đó chính là \textbf{Keyword Spotting} (KWS) --
nhận dạng từ khóa thoại.

Dự án của chúng ta xây dựng một hệ thống tương tự, nhưng có \textbf{ba điểm đặc biệt}:

\begin{enumerate}[leftmargin=2cm]
    \item \textbf{Học nhanh (Few-Shot):}
    Các hệ thống thông thường cần hàng nghìn mẫu giọng nói để học một từ khóa mới.
    Hệ thống của chúng ta chỉ cần \textbf{3 đến 5 mẫu}. Giống như bạn chỉ cần nghe
    tên một người vài lần là nhớ được, không cần nghe hàng nghìn lần.

    \item \textbf{Biết nói ``Tôi không biết'' (Open-Set):}
    Nếu ai đó nói một từ mà hệ thống chưa từng được dạy, thay vì đoán bừa,
    hệ thống sẽ trả lời: ``\textit{Đây là từ lạ, tôi không biết}.'' Điều này rất
    quan trọng trong thực tế -- bạn không muốn loa thông minh bật nhạc khi bạn
    chỉ đang nói chuyện bình thường.

    \item \textbf{Hoạt động trong môi trường ồn:}
    Hệ thống được huấn luyện để nghe được từ khóa ngay cả khi xung quanh có
    tiếng ồn -- quán cà phê, đường phố, văn phòng. Giống như bạn vẫn nghe được
    tên mình dù đang ở nơi đông đúc.
\end{enumerate}

\begin{examplebox}
\textbf{Ví dụ thực tế:} Bạn có một loa thông minh ở nhà. Bạn dạy nó từ khóa
``bật đèn'' bằng cách nói từ đó 5 lần. Từ đó về sau, mỗi khi bạn nói ``bật đèn'',
hệ thống nhận ra và thực hiện lệnh. Nhưng nếu bạn nói ``bật quạt'' (chưa được dạy),
hệ thống sẽ nói: ``Tôi chưa biết từ này'' thay vì bật đèn.
\end{examplebox}


% ============================================================
\section{Âm thanh được xử lý như thế nào?}
% ============================================================

Máy tính không thể ``nghe'' giọng nói giống con người. Nó chỉ hiểu các con số.
Vì vậy, bước đầu tiên là \textbf{chuyển âm thanh thành các con số} mà máy tính
có thể làm việc được.

\subsection{Từ sóng âm đến ``ảnh chụp'' âm thanh}

Khi bạn nói một từ, micro thu lại sóng âm -- đó là một đường cong lên xuống
liên tục. Nhưng đường cong này quá phức tạp để máy tính hiểu trực tiếp.

Chúng ta dùng một kỹ thuật gọi là \textbf{MFCC} để ``chụp ảnh'' âm thanh.
Quá trình này giống như \textbf{chuyển một bài hát thành bản nhạc (sheet music)}:
bài hát thì liên tục và phức tạp, nhưng bản nhạc ghi lại các nốt nhạc rời rạc,
dễ đọc hơn rất nhiều.

\begin{center}
\begin{tikzpicture}[
    block/.style={rectangle, draw, rounded corners, minimum height=1.2cm,
                  minimum width=3cm, align=center, font=\small},
    arrow/.style={-{Stealth[length=3mm]}, thick}
]
    \node[block, fill=core!15] (wav) {Sóng âm thanh\\(giọng nói thô)};
    \node[block, fill=core!25, right=1.5cm of wav] (mfcc) {Tính MFCC\\(40 đặc trưng)};
    \node[block, fill=core!35, right=1.5cm of mfcc] (narrow) {Giữ 10 đặc trưng\\quan trọng nhất};
    \node[block, fill=core!45, right=1.5cm of narrow, text=white] (out) {``Ảnh chụp''\\âm thanh};

    \draw[arrow] (wav) -- (mfcc);
    \draw[arrow] (mfcc) -- (narrow);
    \draw[arrow] (narrow) -- (out);

    \node[below=0.3cm of wav, font=\scriptsize, text=gray] {1 giây, 16.000 mẫu};
    \node[below=0.3cm of mfcc, font=\scriptsize, text=gray] {40 nét, 49 khung};
    \node[below=0.3cm of narrow, font=\scriptsize, text=gray] {10 nét, 49 khung};
    \node[below=0.3cm of out, font=\scriptsize, text=gray] {Ma trận số};
\end{tikzpicture}
\end{center}

\subsection{Tại sao chỉ giữ 10 đặc trưng?}

Kỹ thuật MFCC ban đầu tính ra 40 ``nét đặc trưng'' cho mỗi khung thời gian.
Nhưng giống như khi bạn mô tả một người cho ai đó, bạn chỉ cần nói vài đặc
điểm nổi bật (cao, tóc ngắn, đeo kính) chứ không cần mô tả từng sợi tóc.

Chúng ta \textbf{chỉ giữ 10 đặc trưng đầu tiên} vì chúng chứa hầu hết thông
tin quan trọng. 30 đặc trưng còn lại chủ yếu là chi tiết thừa, bỏ đi giúp hệ
thống chạy nhanh hơn mà không mất độ chính xác.

\begin{examplebox}
\textbf{So sánh dễ hiểu:} 1 giây âm thanh được chuyển thành một bảng số có
49 cột (49 ``lát cắt'' thời gian) và 10 hàng (10 đặc trưng). Tổng cộng 490
con số -- đủ để máy tính ``hiểu'' bạn vừa nói gì.
\end{examplebox}


% ============================================================
\section{Bộ não của hệ thống -- Mạng DSCNN}
% ============================================================

Sau khi có ``ảnh chụp'' âm thanh (490 con số), chúng ta cần một \textbf{bộ não}
để hiểu ý nghĩa của nó. Bộ não này là một \textbf{mạng nơ-ron} (neural network)
có tên DSCNN.

\subsection{Mạng nơ-ron là gì?}

Hãy tưởng tượng bạn đang dạy một đứa trẻ phân biệt các loại quả:

\begin{itemize}[leftmargin=1.5cm]
    \item Lần đầu, đứa trẻ nhìn thấy quả táo đỏ, tròn. Nó ghi nhớ: đỏ + tròn = táo.
    \item Lần sau, nó thấy quả táo xanh. Nó học thêm: à, táo không nhất thiết phải đỏ.
    \item Sau nhiều lần, nó biết nhận ra táo dựa trên hình dáng, kích thước, cuống quả...
          chứ không chỉ dựa vào màu sắc.
\end{itemize}

Mạng nơ-ron hoạt động tương tự. Nó \textbf{học từ rất nhiều ví dụ} để tìm ra
các đặc điểm quan trọng nhất giúp phân biệt các từ khóa.

\subsection{DSCNN -- nhẹ và nhanh}

DSCNN là viết tắt của ``Depthwise Separable CNN'' -- nhưng bạn không cần nhớ
tên dài dòng đó. Điều quan trọng là:

\begin{itemize}[leftmargin=1.5cm]
    \item Nó \textbf{rất nhẹ} -- có thể chạy trên điện thoại hoặc thiết bị nhỏ.
    \item Nó \textbf{rất nhanh} -- xử lý 1 giây âm thanh trong chưa đầy 10 mili-giây
          (nhanh hơn một cái chớp mắt).
    \item Nó biến ``ảnh chụp'' âm thanh (490 con số) thành \textbf{một điểm
          trong không gian} (276 con số). Điểm này gọi là \textbf{embedding}.
\end{itemize}

\begin{examplebox}
\textbf{Hình dung đơn giản:} Tưởng tượng bạn có rất nhiều bản ghi âm khác nhau.
DSCNN giống như một cỗ máy ``ép'' mỗi bản ghi thành \textbf{một chấm nhỏ trên bản đồ}.
Các bản ghi cùng từ khóa sẽ có chấm gần nhau, các bản ghi khác từ khóa sẽ có
chấm xa nhau.
\end{examplebox}

\begin{center}
\begin{tikzpicture}[scale=0.9]
    % Truoc khi qua DSCNN
    \node[font=\small\bfseries] at (-4, 3) {Trước DSCNN};
    \node[font=\small\bfseries] at (4, 3) {Sau DSCNN};

    % Left side - raw audio
    \draw[thick, gray, rounded corners] (-6.5, -2.5) rectangle (-1.5, 2.5);
    \node[font=\scriptsize, text=gray] at (-4, 2) {490 con số -- rối rắm};
    \foreach \i in {1,...,15}{
        \pgfmathsetmacro{\x}{-6.2 + rand*4.2}
        \pgfmathsetmacro{\y}{-2.0 + rand*3.5}
        \fill[gray!60] (\x, \y) circle (2pt);
    }

    % Arrow
    \draw[-{Stealth[length=4mm]}, very thick, core] (-1, 0) -- (1, 0)
        node[midway, above, font=\small] {DSCNN};

    % Right side - embeddings
    \draw[thick, gray, rounded corners] (1.5, -2.5) rectangle (6.5, 2.5);
    \node[font=\scriptsize, text=gray] at (4, 2) {276 số -- có trật tự};

    % Cluster "yes"
    \fill[core] (2.5, 1.0) circle (3pt);
    \fill[core] (2.8, 0.7) circle (3pt);
    \fill[core] (2.6, 0.5) circle (3pt);
    \node[font=\scriptsize, core] at (2.6, 1.5) {``yes''};

    % Cluster "stop"
    \fill[ext1] (5.0, 1.2) circle (3pt);
    \fill[ext1] (5.3, 0.9) circle (3pt);
    \fill[ext1] (5.1, 0.7) circle (3pt);
    \node[font=\scriptsize, ext1] at (5.1, 1.7) {``stop''};

    % Cluster "go"
    \fill[ext2] (3.8, -1.5) circle (3pt);
    \fill[ext2] (4.1, -1.2) circle (3pt);
    \fill[ext2] (3.9, -1.8) circle (3pt);
    \node[font=\scriptsize, ext2] at (3.9, -0.7) {``go''};
\end{tikzpicture}
\end{center}


% ============================================================
\section{Học từ 3--5 mẫu -- Few-Shot Learning}
% ============================================================

Đây là phần thú vị nhất của dự án. Hệ thống thông thường cần hàng nghìn mẫu
để học một từ mới. Hệ thống của chúng ta chỉ cần \textbf{3 đến 5 mẫu}.

\subsection{Cách hoạt động}

Giả sử bạn muốn dạy hệ thống nhận ra từ ``bật đèn'':

\begin{enumerate}[leftmargin=2cm]
    \item Bạn nói ``bật đèn'' \textbf{5 lần}. Mỗi lần, hệ thống thu âm và
          chuyển thành một ``chấm'' trên bản đồ (thông qua MFCC + DSCNN).
    \item Hệ thống tính \textbf{điểm trung tâm} của 5 chấm đó. Điểm trung tâm
          này gọi là \textbf{prototype} -- ``đại diện'' cho từ ``bật đèn''.
    \item Khi có ai nói một từ mới, hệ thống chuyển nó thành chấm, rồi xem
          chấm đó \textbf{gần prototype nào nhất}.
\end{enumerate}

\begin{center}
\begin{tikzpicture}[scale=1.1]
    % 5 samples
    \fill[core!70] (0.0, 0.3) circle (3pt);
    \fill[core!70] (0.4, -0.2) circle (3pt);
    \fill[core!70] (-0.3, -0.1) circle (3pt);
    \fill[core!70] (0.2, 0.5) circle (3pt);
    \fill[core!70] (-0.1, -0.4) circle (3pt);

    % Labels
    \node[font=\scriptsize, core!70, right] at (0.5, 0.3) {mẫu 1};
    \node[font=\scriptsize, core!70, right] at (0.5, -0.2) {mẫu 2};

    % Prototype (center)
    \fill[core, thick] (0.04, 0.02) circle (5pt);
    \node[font=\small\bfseries, core, below right] at (0.2, -0.1) {Prototype};

    % Arrow from samples to prototype
    \draw[dashed, gray, thin] (0.0, 0.3) -- (0.04, 0.02);
    \draw[dashed, gray, thin] (0.4, -0.2) -- (0.04, 0.02);
    \draw[dashed, gray, thin] (-0.3, -0.1) -- (0.04, 0.02);
    \draw[dashed, gray, thin] (0.2, 0.5) -- (0.04, 0.02);
    \draw[dashed, gray, thin] (-0.1, -0.4) -- (0.04, 0.02);

    % Annotation
    \draw[decorate, decoration={brace, amplitude=8pt, mirror}, thick]
        (-0.8, -0.7) -- (-0.8, 0.7) node[midway, left=10pt, font=\small, align=right]
        {5 mẫu\\``bật đèn''};
    \node[font=\small, align=center] at (2.5, -0.8) {Prototype = trung bình\\của 5 mẫu};
\end{tikzpicture}
\end{center}

\begin{examplebox}
\textbf{So sánh:} Hãy nghĩ về ``khuôn mặt trung bình''. Nếu bạn chụp 5 tấm
ảnh của cùng một người ở các góc khác nhau rồi ``trộn'' lại, bạn sẽ có một
hình ảnh đại diện cho khuôn mặt người đó. Prototype hoạt động tương tự --
nó là ``giọng nói trung bình'' của một từ khóa.
\end{examplebox}

\subsection{Quá trình huấn luyện}

Trước khi hệ thống có thể học nhanh từ 3--5 mẫu, bộ não DSCNN cần được
\textbf{huấn luyện trước} trên một tập dữ liệu lớn (khoảng 450 từ tiếng Anh,
mỗi từ hàng nghìn mẫu). Quá trình này giống như:

\begin{examplebox}
\textbf{Ví dụ:} Trước khi bạn có thể nhận ra một giống chim mới chỉ sau vài lần
nhìn, bạn đã có \textbf{nhiều năm kinh nghiệm nhìn các loại chim khác nhau}.
Não bạn đã biết cách chú ý vào mỏ, cánh, màu lông... DSCNN cũng vậy -- nó
``nhìn'' hàng trăm từ khác nhau để học cách tập trung vào các đặc trưng
giọng nói quan trọng.
\end{examplebox}


% ============================================================
\section{Nhận diện và từ chối -- Open-Set}
% ============================================================

Sau khi có prototype cho mỗi từ khóa, việc nhận diện rất đơn giản:

\subsection{Đo khoảng cách}

Khi ai đó nói một từ, hệ thống chuyển nó thành ``chấm'' (embedding), rồi
\textbf{đo khoảng cách} từ chấm đó đến tất cả các prototype.

\begin{itemize}[leftmargin=1.5cm]
    \item \textbf{Gần nhất} = dự đoán đó là từ khóa tương ứng.
          Ví dụ: chấm gần prototype ``bật đèn'' nhất $\rightarrow$ ``bật đèn''.
    \item \textbf{Quá xa tất cả} = ``\textit{Từ lạ, tôi không biết.}''
          Ví dụ: chấm ở rất xa mọi prototype $\rightarrow$ từ chối.
\end{itemize}

\subsection{Vòng tròn chấp nhận}

Chúng ta đặt một ``vòng tròn'' quanh mỗi prototype. Nếu chấm nằm \textbf{trong}
vòng tròn, hệ thống chấp nhận. Nếu nằm \textbf{ngoài} tất cả các vòng tròn,
hệ thống từ chối.

\begin{center}
\begin{tikzpicture}[scale=1.0]
    % Prototype 1: "yes"
    \fill[core] (-2.5, 0.8) circle (5pt);
    \node[font=\small\bfseries, core, above=6pt] at (-2.5, 0.8) {``yes''};
    \draw[core, thick, dashed] (-2.5, 0.8) circle (1.2cm);

    % Prototype 2: "stop"
    \fill[ext1] (2.0, 1.0) circle (5pt);
    \node[font=\small\bfseries, ext1, above=6pt] at (2.0, 1.0) {``stop''};
    \draw[ext1, thick, dashed] (2.0, 1.0) circle (1.2cm);

    % Prototype 3: "go"
    \fill[ext2] (0.0, -1.5) circle (5pt);
    \node[font=\small\bfseries, ext2, above=6pt] at (0.0, -1.5) {``go''};
    \draw[ext2, thick, dashed] (0.0, -1.5) circle (1.2cm);

    % Query 1 - accepted (inside "yes" circle)
    \fill[accepted] (-2.0, 0.3) circle (4pt);
    \node[font=\scriptsize, accepted, right=3pt] at (-2.0, 0.3) {%
        ``yes'' $\checkmark$};

    % Query 2 - rejected (outside all circles)
    \fill[rejected] (0.0, 0.5) circle (4pt);
    \node[font=\scriptsize, rejected, right=3pt] at (0.0, 0.5) {%
        Từ lạ $\times$};

    % Threshold annotation
    \draw[gray, thin, |->] (-2.5, 0.8) -- (-1.3, 0.8);
    \node[font=\scriptsize, gray, above] at (-1.9, 0.8) {bán kính};
\end{tikzpicture}
\end{center}

\begin{examplebox}
\textbf{So sánh:} Hãy tưởng tượng bạn đứng giữa sân, quanh bạn có 3 vòng tròn
được vẽ trên mặt đất, mỗi vòng có một bảng tên (``yes'', ``stop'', ``go'').
Khi ai đó ném một quả bóng vào sân:
\begin{itemize}
    \item Nếu bóng rơi \textbf{trong} vòng ``yes'' $\rightarrow$ ``Đó là từ YES!''
    \item Nếu bóng rơi \textbf{ngoài} mọi vòng $\rightarrow$ ``Không biết từ này.''
\end{itemize}
Bán kính vòng tròn chính là ``ngưỡng chấp nhận'' -- to quá thì dễ nhận nhầm,
nhỏ quá thì hay bỏ sót.
\end{examplebox}

\subsection{Tại sao cách này hay hơn?}

Nhiều hệ thống khác dùng ``xác suất'' để quyết định (ví dụ: 80\% là ``yes'',
15\% là ``stop'', 5\% là ``go''). Nhưng cách đó có vấn đề: hệ thống buộc phải
chọn một đáp án ngay cả khi từ đó hoàn toàn lạ.

Cách ``vòng tròn chấp nhận'' của chúng ta đơn giản hơn và trực quan hơn:
\textbf{gần thì nhận, xa thì từ chối.}


% ============================================================
\section{Chống nhiễu -- Nghe rõ trong môi trường ồn}
% ============================================================

Trong thực tế, hiếm khi bạn nói trong phòng im lặng hoàn toàn. Xung quanh
luôn có tiếng ồn: xe cộ, quạt máy, người nói chuyện... Hệ thống cần phải
nghe được từ khóa \textbf{ngay cả khi có nhiễu}.

\subsection{Khi huấn luyện: ``tập nghe'' trong ồn ào}

Khi dạy bộ não DSCNN, chúng ta \textbf{cố ý trộn thêm tiếng ồn} vào mẫu
âm thanh sạch. Như vậy, DSCNN ``quen'' với việc nghe trong môi trường ồn
ngay từ khi còn đang học.

\begin{center}
\begin{tikzpicture}[
    block/.style={rectangle, draw, rounded corners, minimum height=1cm,
                  minimum width=2.5cm, align=center, font=\small},
    arrow/.style={-{Stealth[length=3mm]}, thick}
]
    \node[block, fill=accepted!15] (clean) {Giọng nói sạch\\``bật đèn''};
    \node[block, fill=rejected!15, below=0.8cm of clean] (noise) {Tiếng ồn\\(quán cà phê)};
    \node[font=\Huge] at (3.5, -0.5) {+};
    \node[block, fill=ext2!20, right=3cm of clean, yshift=-0.5cm] (mixed) {Giọng nói\\+ nhiễu};

    \draw[arrow] (clean.east) -- ++(0.5,0) |- (mixed.west);
    \draw[arrow] (noise.east) -- ++(0.5,0) |- (mixed.west);
\end{tikzpicture}
\end{center}

\begin{examplebox}
\textbf{Ví dụ:} Giống như sinh viên học tiếng Anh trong phòng im lặng sẽ gặp
khó khăn khi giao tiếp ở sân bay ồn ào. Nhưng nếu \textbf{luyện tập nghe trong
quán cà phê} ngay từ đầu, họ sẽ quen và nghe tốt hơn trong mọi tình huống.
\end{examplebox}

\subsection{Khi sử dụng: lọc nhiễu trước khi nhận diện}

Ngoài việc ``tập nghe'' trong ồn ào, hệ thống còn có thêm một bước \textbf{lọc nhiễu}
trước khi phân tích âm thanh. Bước này giống như đeo tai nghe chống ồn --
loại bỏ phần lớn tiếng ồn nền trước khi bạn nghe.


% ============================================================
\section{Nghe liên tục -- Streaming}
% ============================================================

Các phần trước giả định bạn đưa cho hệ thống \textbf{một file âm thanh ngắn}
(1 giây) để phân tích. Nhưng trong thực tế, hệ thống phải \textbf{lắng nghe
liên tục} và phát hiện từ khóa bất kỳ lúc nào.

\subsection{Bước 1: Kiểm tra ``có ai nói không?'' (VAD)}

Hệ thống không chạy nhận diện từ khóa 100\% thời gian -- điều đó quá tốn năng lượng.
Thay vào đó, nó có một ``bảo vệ'' đứng ở cổng, gọi là \textbf{VAD}
(Voice Activity Detection):

\begin{itemize}[leftmargin=1.5cm]
    \item Nếu \textbf{không có giọng nói} (im lặng hoặc chỉ có tiếng ồn nền):
          không làm gì cả, tiết kiệm pin.
    \item Nếu \textbf{có giọng nói}: ``Có người nói! Kích hoạt nhận diện!''
\end{itemize}

\begin{examplebox}
\textbf{Ví dụ:} Giống bảo vệ ở cổng công ty. Bảo vệ chỉ kiểm tra thẻ khi
\textbf{có người đến}. Nếu không ai đến, bảo vệ không cần làm gì. Nhờ vậy,
hệ thống tiết kiệm rất nhiều năng lượng.
\end{examplebox}

\subsection{Bước 2: Cửa sổ trượt}

Khi VAD phát hiện giọng nói, hệ thống cắt ra \textbf{1 giây âm thanh} để
phân tích. Sau đó, cửa sổ ``trượt'' về phía trước 0,5 giây và cắt tiếp
1 giây nữa. Quá trình lặp lại liên tục.

\begin{center}
\begin{tikzpicture}[scale=0.85]
    % Timeline
    \draw[thick, gray] (0, 0) -- (12, 0);
    \foreach \x in {0, 1, ..., 12} {
        \draw[gray] (\x, -0.1) -- (\x, 0.1);
    }
    \node[font=\scriptsize, gray, below] at (0, -0.15) {0s};
    \node[font=\scriptsize, gray, below] at (4, -0.15) {2s};
    \node[font=\scriptsize, gray, below] at (8, -0.15) {4s};
    \node[font=\scriptsize, gray, below] at (12, -0.15) {6s};

    % Window 1
    \draw[very thick, core, rounded corners=2pt] (0, 0.3) rectangle (2, 0.8);
    \node[font=\scriptsize, core] at (1, 0.55) {Cửa sổ 1};

    % Window 2
    \draw[very thick, ext1, rounded corners=2pt] (1, 1.0) rectangle (3, 1.5);
    \node[font=\scriptsize, ext1] at (2, 1.25) {Cửa sổ 2};

    % Window 3
    \draw[very thick, ext2, rounded corners=2pt] (2, 1.7) rectangle (4, 2.2);
    \node[font=\scriptsize, ext2] at (3, 1.95) {Cửa sổ 3};

    % Annotation
    \draw[{Stealth[length=2mm]}-{Stealth[length=2mm]}, thick] (0, -0.6) -- (2, -0.6)
        node[midway, below, font=\scriptsize] {1 giây};
    \draw[-{Stealth[length=2mm]}, thick, gray] (0.5, -1.1) -- (1.5, -1.1)
        node[midway, below, font=\scriptsize] {trượt 0,5s};
\end{tikzpicture}
\end{center}

Mỗi cửa sổ 1 giây được đưa qua toàn bộ quy trình: lọc nhiễu $\rightarrow$
MFCC $\rightarrow$ DSCNN $\rightarrow$ so sánh prototype $\rightarrow$ kết quả.
Toàn bộ quá trình này hoàn thành trong \textbf{dưới 100 mili-giây} --
nhanh hơn một cái chớp mắt, nên người dùng không cảm thấy độ trễ.


% ============================================================
\section{Tổng kết -- Toàn bộ quy trình}
% ============================================================

Bây giờ, hãy ghép tất cả lại thành bức tranh toàn cảnh:

\begin{center}
\begin{tikzpicture}[
    step/.style={rectangle, draw, rounded corners=5pt, minimum height=1.4cm,
                 minimum width=2.8cm, align=center, font=\small, thick},
    arrow/.style={-{Stealth[length=3mm]}, very thick},
    note/.style={font=\scriptsize, text=gray, align=center}
]
    % Steps
    \node[step, fill=ext2!20] (s1) at (0, 0) {\textbf{1. Nghe}\\Thu âm\\liên tục};
    \node[step, fill=ext2!30] (s2) at (3.5, 0) {\textbf{2. Lọc}\\Có giọng nói?\\(VAD)};
    \node[step, fill=ext1!20] (s3) at (7, 0) {\textbf{3. Sạch}\\Lọc nhiễu\\nền};
    \node[step, fill=core!20] (s4) at (0, -3) {\textbf{4. Chụp}\\MFCC\\(ảnh âm thanh)};
    \node[step, fill=core!30] (s5) at (3.5, -3) {\textbf{5. Hiểu}\\DSCNN\\(embedding)};
    \node[step, fill=core!40, text=white] (s6) at (7, -3) {\textbf{6. Quyết định}\\Gần prototype?\\Nhận / Từ chối};

    % Arrows
    \draw[arrow, ext2] (s1) -- (s2);
    \draw[arrow, ext2] (s2) -- (s3) node[midway, above, font=\scriptsize] {có};
    \draw[arrow, ext1] (s3) -- (s3 |- s4.north) -- (s4);
    \draw[arrow, core] (s4) -- (s5);
    \draw[arrow, core] (s5) -- (s6);

    % Reject path from VAD
    \node[note, below=0.2cm of s2] {không $\rightarrow$ bỏ qua};

    % Output
    \node[step, fill=accepted!20, right=1cm of s6] (out) {\textbf{Kết quả}\\``bật đèn''\\hoặc ``?''};
    \draw[arrow] (s6) -- (out);

    % Notes
    \node[note, below=0.15cm of s1] {micro / file};
    \node[note, below=0.15cm of s4] {490 con số};
    \node[note, below=0.15cm of s5] {276 con số};
\end{tikzpicture}
\end{center}

\vspace{0.5cm}

\begin{center}
\begin{tabular}{clll}
\toprule
\textbf{Bước} & \textbf{Tên} & \textbf{Làm gì} & \textbf{So sánh} \\
\midrule
1 & Nghe & Thu âm liên tục & Micro bật sẵn \\
2 & Lọc giọng & Kiểm tra có người nói? & Bảo vệ ở cổng \\
3 & Lọc nhiễu & Bỏ tiếng ồn nền & Tai nghe chống ồn \\
4 & Chụp ảnh & Chuyển sóng âm thành số & Chụp ảnh bài hát \\
5 & Hiểu & Tìm ý nghĩa trong số & Não nhận ra ý nghĩa \\
6 & Quyết định & Gần từ nào? Hay từ lạ? & Bóng rơi trong vòng? \\
\bottomrule
\end{tabular}
\end{center}

\vspace{1cm}

\begin{examplebox}
\textbf{Tóm tắt bằng một câu:}
Hệ thống lắng nghe liên tục, lọc bỏ tiếng ồn, chuyển giọng nói thành con số,
rồi so sánh với các ``mẫu đại diện'' đã học trước đó. Nếu giống -- nhận ra.
Nếu khác hết -- từ chối. Tất cả diễn ra trong chưa đầy 1/10 giây.
\end{examplebox}


\end{document}
